% ============================================================
%  LAMPIRAN 7 – DAFTAR BUKU TEKS / KORPUS PENELITIAN
%  Compile standalone:  pdflatex lampiran-7.tex
%  Compile as part of thesis: pdflatex main.tex
% ============================================================
\documentclass[../../thesis]{subfiles}
\usepackage{hyphenat}
\usepackage{iftex}

% ── Encoding, Language & Fonts ───────────────────────────────
\ifXeTeX
    \usepackage{polyglossia}
    \setmainlanguage{bahasai}
    \usepackage{fontspec}
    \setmainfont{TeX Gyre Termes}
    \setsansfont{TeX Gyre Heros}
    \setmonofont[Scale=0.88]{TeX Gyre Cursor}
\else
    \usepackage[utf8]{inputenc}
    \usepackage[T1]{fontenc}
    \usepackage[indonesian]{babel}
    \usepackage{mathptmx}
    \usepackage{helvet}
    \usepackage{courier}
\fi

% ── Page geometry (A4, UI thesis margins) ────────────────────
\usepackage[
  a4paper,
  headheight = 0pt,
  top    = 3cm,
  bottom = 3cm,
  left   = 3cm,
  right  = 3cm
]{geometry}

% ── Line spacing & paragraph ─────────────────────────────────
\usepackage{setspace}
\onehalfspacing
\setlength{\parindent}{1.27cm}
\setlength{\parskip}{0pt}

% ── Microtype & justification ────────────────────────────────
\usepackage{microtype}
\sloppy

% ── Headings ─────────────────────────────────────────────────
\usepackage{titlesec}

\titleformat{\chapter}[block]
  {\centering\rmfamily\bfseries\large}{}
  {0pt}{\MakeUppercase}
\titlespacing*{\chapter}{0pt}{0pt}{1.5em}

\titleformat{\section}
  {\rmfamily\bfseries\normalsize}{\thesection}{1em}{}
\titlespacing*{\section}{0pt}{1.2em}{0.6em}

\titleformat{\subsection}
  {\rmfamily\bfseries\normalsize}{\thesubsection}{1em}{}
\titlespacing*{\subsection}{0pt}{1em}{0.5em}

% ── Tables ───────────────────────────────────────────────────
\usepackage{longtable, booktabs, array, multirow}
\usepackage{calc}
\usepackage{tabularx}

% ── Hyperlinks ───────────────────────────────────────────────
\usepackage[hidelinks]{hyperref}

% Bibliography setup is inherited from main.tex (biblatex + references.bib)
% when this file is compiled standalone via the subfiles class.

% ============================================================
\begin{document}

% When standalone: switch to appendix mode and set the counter so the
% appendix heading renders correctly. When compiled inside main.tex,
% main.tex issues \appendix and advances the counter before \subfile.
\ifSubfilesClassLoaded{}{\appendix\setcounter{chapter}{1}}

\lampiran{Daftar Buku Teks (Korpus Penelitian)}
\label{lamp:korpus}

Lampiran ini merinci korpus buku teks yang dipakai pada Tahap~1--2
\textit{pipeline} (Subbab~\ref{subsec:data-primer}). Korpus primer adalah tiga
\textbf{Buku Siswa} Fisika, Kimia, dan Biologi Fase~F Kelas~XII yang menjadi
sumber ekstraksi entitas dan relasi. Tiga \textbf{Buku Panduan Guru} pasangannya
disertakan sebagai sumber pendukung, khusus untuk validasi ``Materi Pokok'' pada
tahap ekstraksi (Subbab~\ref{sec:ingestion}, butir Ekstraksi materi pokok).
Seluruh buku merupakan terbitan resmi Pusat Perbukuan, Badan Standar, Kurikulum,
dan Asesmen Pendidikan, Kementerian Pendidikan, Kebudayaan, Riset, dan Teknologi
(Kemendikbudristek), cetakan pertama 2022, dan diperoleh dari repositori resmi
\url{https://buku.kemdikbud.go.id}.

\begin{longtable}{@{}
  >{\raggedright\arraybackslash}p{0.06\linewidth}
  >{\raggedright\arraybackslash}p{0.30\linewidth}
  >{\raggedright\arraybackslash}p{0.42\linewidth}
  >{\raggedright\arraybackslash}p{0.16\linewidth}@{}}
\caption{Korpus buku teks Kemendikbudristek Fase~F Kelas~XII (cetakan pertama,
  2022).}
\label{tab:korpus}\\
\toprule
\textbf{No.} & \textbf{Judul} & \textbf{Penulis} & \textbf{ISBN (jil.~2)} \\
\midrule
\endfirsthead
\toprule
\textbf{No.} & \textbf{Judul} & \textbf{Penulis} & \textbf{ISBN (jil.~2)} \\
\midrule
\endhead
\midrule
\multicolumn{4}{r@{}}{\small\itshape bersambung ke halaman berikutnya}\\
\endfoot
\bottomrule
\endlastfoot
\multicolumn{4}{@{}l}{\textit{Buku Siswa (korpus primer ekstraksi)}}\\
\midrule
1 & Fisika untuk SMA/MA Kelas XII
  & Lia Laela Sarah dan Irma Rahma Suwarma
  & 978-623-472-722-7 \\
2 & Kimia untuk SMA/MA Kelas XII
  & Galuh Yuliani, Hanhan Dianhar, dan Aang Suhendar
  & 978-602-427-968-4 \\
3 & Biologi untuk SMA/MA Kelas XII
  & Shilviani Dewi, Amalia Shari, Rani Elisa Purba, dan Remigius Gunawan
    Susilowarno
  & 978-602-427-958-5 \\
\midrule
\multicolumn{4}{@{}l}{\textit{Buku Panduan Guru (sumber pendukung ``Materi Pokok'')}}\\
\midrule
4 & Buku Panduan Guru Fisika untuk SMA/MA Kelas XII
  & Lia Laela Sarah dan Irma Rahma Suwarma
  & 978-623-472-725-8 \\
5 & Buku Panduan Guru Kimia untuk SMA/MA Kelas XII
  & Galuh Yuliani, Hanhan Dianhar, dan Aang Suhendar
  & 978-602-427-969-1 \\
6 & Buku Panduan Guru Biologi untuk SMA/MA Kelas XII
  & Shilviani Dewi, Amalia Shari, Rani Elisa Purba, dan Remigius Gunawan
    Susilowarno
  & 978-602-427-959-2 \\
\end{longtable}

\noindent\textit{Catatan.} ISBN yang dicantumkan adalah nomor jilid~2 (Kelas~XII)
dari masing-masing judul; setiap judul juga memiliki ISBN nomor jilid lengkap.

\end{document}
